% ============================================================================
% DOCUMENTACIÓN TÉCNICA - PROYECTO SAFEAREA
% Formato: LaTeX
% Fecha: 27 de Noviembre de 2025
% ============================================================================

\documentclass[12pt,a4paper]{article}

% Paquetes necesarios
\usepackage[utf8]{inputenc}
\usepackage[spanish]{babel}
\usepackage{geometry}
\usepackage{graphicx}
\usepackage{hyperref}
\usepackage{listings}
\usepackage{xcolor}
\usepackage{booktabs}
\usepackage{longtable}
\usepackage{array}
\usepackage{fancyhdr}
\usepackage{titlesec}
\usepackage{tocloft}

% Configuración de página
\geometry{margin=2.5cm}
\setlength{\parindent}{0pt}
\setlength{\parskip}{0.5em}

% Configuración de colores
\definecolor{codegreen}{rgb}{0,0.6,0}
\definecolor{codegray}{rgb}{0.5,0.5,0.5}
\definecolor{codepurple}{rgb}{0.58,0,0.82}
\definecolor{backcolour}{rgb}{0.95,0.95,0.92}
\definecolor{primaryblue}{RGB}{74,105,255}

% Configuración de código
\lstdefinestyle{mystyle}{
    backgroundcolor=\color{backcolour},   
    commentstyle=\color{codegreen},
    keywordstyle=\color{magenta},
    numberstyle=\tiny\color{codegray},
    stringstyle=\color{codepurple},
    basicstyle=\ttfamily\footnotesize,
    breakatwhitespace=false,         
    breaklines=true,                 
    captionpos=b,                    
    keepspaces=true,                 
    numbers=left,                    
    numbersep=5pt,                  
    showspaces=false,                
    showstringspaces=false,
    showtabs=false,                  
    tabsize=2
}
\lstset{style=mystyle}

% Configuración de encabezados
\pagestyle{fancy}
\fancyhf{}
\rhead{SafeArea - Documentación Técnica}
\lhead{v0.1.0}
\rfoot{Página \thepage}

% Configuración de hipervínculos
\hypersetup{
    colorlinks=true,
    linkcolor=primaryblue,
    filecolor=magenta,      
    urlcolor=cyan,
    pdftitle={SafeArea - Documentación Técnica},
}

% ============================================================================
% INICIO DEL DOCUMENTO
% ============================================================================

\begin{document}

% Portada
\begin{titlepage}
    \centering
    \vspace*{2cm}
    
    {\Huge\bfseries SAFEAREA\par}
    \vspace{0.5cm}
    {\Large Aplicación de Seguridad Comunitaria\par}
    \vspace{2cm}
    
    {\LARGE\bfseries Documentación Técnica\par}
    \vspace{1cm}
    
    {\large Versión 0.1.0\par}
    \vspace{2cm}
    
    \begin{tabular}{ll}
        \textbf{Plataforma:} & Flutter (Android, iOS, Web) \\
        \textbf{Backend:} & Firebase \\
        \textbf{Base de datos:} & Cloud Firestore \\
        \textbf{Fecha:} & 27 de Noviembre de 2025 \\
    \end{tabular}
    
    \vfill
    
    {\large com.newsdevs.safearea\par}
\end{titlepage}

% Tabla de contenidos
\tableofcontents
\newpage

% ============================================================================
% SECCIÓN 1: RESUMEN EJECUTIVO
% ============================================================================

\section{Resumen Ejecutivo}

\subsection{Descripción del Proyecto}

\textbf{SafeArea} es una aplicación móvil de \textbf{seguridad comunitaria} desarrollada en Flutter que permite a los ciudadanos reportar incidentes de seguridad en tiempo real, comunicarse mediante chat grupal por zonas geográficas y recibir notificaciones sobre eventos relevantes en su comunidad.

\subsection{Problema que Resuelve}

\begin{itemize}
    \item \textbf{Falta de comunicación ciudadana:} Permite reportar incidentes (robos, incendios, emergencias, accidentes) de forma inmediata.
    \item \textbf{Descoordinación comunitaria:} Facilita la comunicación entre vecinos mediante chats grupales organizados por zonas (Centro, Sur, Norte de Tacna).
    \item \textbf{Respuesta lenta ante emergencias:} Sistema de notificaciones push en tiempo real para alertar a la comunidad.
    \item \textbf{Gestión de seguridad:} Panel de administración para moderar reportes y gestionar usuarios.
\end{itemize}

\subsection{Tecnologías Principales}

\begin{table}[h]
\centering
\begin{tabular}{|l|l|}
\hline
\textbf{Categoría} & \textbf{Tecnología} \\
\hline
Framework & Flutter 3.7.2+ (Dart) \\
Backend & Firebase (BaaS) \\
Base de datos & Cloud Firestore \\
Autenticación & Firebase Auth \\
Almacenamiento & Firebase Storage \\
Notificaciones & Firebase Cloud Messaging \\
Mapas & Google Maps Flutter + Flutter Map (OSM) \\
Estado & Provider (ChangeNotifier) \\
\hline
\end{tabular}
\caption{Stack tecnológico principal}
\end{table}

% ============================================================================
% SECCIÓN 2: STACK TECNOLÓGICO
% ============================================================================

\section{Stack Tecnológico}

\subsection{Lenguajes y Frameworks}

\begin{table}[h]
\centering
\begin{tabular}{|l|l|l|}
\hline
\textbf{Componente} & \textbf{Versión} & \textbf{Propósito} \\
\hline
Dart SDK & \^{}3.7.2 & Lenguaje de programación \\
Flutter SDK & 3.x & Framework de UI multiplataforma \\
\hline
\end{tabular}
\caption{Lenguajes y frameworks}
\end{table}

\subsection{Dependencias de Producción}

\begin{lstlisting}[language=yaml, caption=pubspec.yaml - Dependencias]
dependencies:
  # Core Flutter
  flutter: sdk
  flutter_localizations: sdk
  
  # Firebase Suite
  firebase_core: ^4.2.0
  firebase_auth: ^6.1.1
  cloud_firestore: ^6.0.3
  firebase_storage: ^13.0.3
  firebase_messaging: ^16.0.3
  
  # Estado y Persistencia
  provider: ^6.1.5+1
  shared_preferences: ^2.5.3
  
  # Mapas y Geolocalizacion
  google_maps_flutter: ^2.9.0
  google_maps_flutter_web: ^0.5.7
  flutter_map: ^7.0.2
  latlong2: ^0.9.1
  geolocator: ^13.0.1
  
  # UI y Media
  image_picker: ^1.1.2
  flutter_local_notifications: ^18.0.1
  http: ^1.2.2
  flutter_launcher_icons: ^0.14.4
\end{lstlisting}

\subsection{Base de Datos (Cloud Firestore)}

\begin{table}[h]
\centering
\begin{tabular}{|l|l|}
\hline
\textbf{Colección} & \textbf{Descripción} \\
\hline
users & Usuarios registrados (perfil, rol, estado) \\
reports & Reportes de incidentes \\
chatGroups & Grupos de chat (zonas y privados) \\
chatGroups/\{id\}/messages & Mensajes de cada grupo \\
\hline
\end{tabular}
\caption{Colecciones de Firestore}
\end{table}

% ============================================================================
% SECCIÓN 3: ARQUITECTURA DEL SOFTWARE
% ============================================================================

\section{Arquitectura del Software}

\subsection{Patrón de Diseño}

El proyecto implementa una arquitectura \textbf{MVVM simplificada} (Model-View-ViewModel) usando \textbf{Provider} como gestor de estado.

\begin{lstlisting}[caption=Estructura de capas]
PRESENTATION (Screens, Widgets, Theme)
         |
         v
STATE MANAGEMENT (AuthService, ReportService, ChatService)
         |
         v
DATA (UserModel, ReportModel, ChatGroupModel)
         |
         v
FIREBASE (Firestore, Auth, Storage)
\end{lstlisting}

\subsection{Estructura de Carpetas}

\begin{lstlisting}[caption=Estructura del proyecto]
lib/
|-- main.dart                 # Punto de entrada
|-- firebase_options.dart     # Configuracion Firebase
|
|-- models/                   # Modelos de datos
|   |-- user_model.dart
|   |-- report_model.dart
|   |-- chat_group_model.dart
|
|-- services/                 # Logica de negocio
|   |-- auth_service.dart
|   |-- report_service.dart
|   |-- chat_service.dart
|   |-- notification_service.dart
|   |-- storage_service.dart
|   |-- theme_service.dart
|   |-- language_service.dart
|
|-- screens/                  # Vistas (UI)
|   |-- login_screen.dart
|   |-- register_screen.dart
|   |-- home_screen.dart
|   |-- reports_screen.dart
|   |-- add_report_screen.dart
|   |-- chat_groups_screen.dart
|   |-- chat_screen.dart
|   |-- profile_screen.dart
|   |-- settings_screen.dart
|   |-- dashboard_screen.dart
|   |-- users_screen.dart
|
|-- widgets/                  # Componentes reutilizables
|   |-- auth_text_field.dart
|   |-- report_card.dart
|   |-- filter_chip.dart
|   |-- loading_button.dart
|
|-- theme/                    # Configuracion de estilos
    |-- app_theme.dart
\end{lstlisting}

% ============================================================================
% SECCIÓN 4: REQUERIMIENTOS FUNCIONALES
% ============================================================================

\section{Especificación de Requerimientos Funcionales (RF)}

\subsection{Módulo de Autenticación}

\begin{longtable}{|p{1.5cm}|p{4cm}|p{8cm}|}
\hline
\textbf{ID} & \textbf{Requerimiento} & \textbf{Descripción} \\
\hline
\endfirsthead
\hline
\textbf{ID} & \textbf{Requerimiento} & \textbf{Descripción} \\
\hline
\endhead
RF-01 & Registro de usuarios & El sistema permite registrar nuevos usuarios con email, contraseña y nombre \\
\hline
RF-02 & Inicio de sesión & El sistema permite autenticar usuarios mediante email y contraseña \\
\hline
RF-03 & Cierre de sesión & El sistema permite cerrar la sesión del usuario actual \\
\hline
RF-04 & Gestión de perfil & El usuario puede ver y editar su nombre, teléfono y foto de perfil \\
\hline
\caption{Requerimientos funcionales - Autenticación}
\end{longtable}

\subsection{Módulo de Reportes}

\begin{longtable}{|p{1.5cm}|p{4cm}|p{8cm}|}
\hline
\textbf{ID} & \textbf{Requerimiento} & \textbf{Descripción} \\
\hline
\endfirsthead
\hline
\textbf{ID} & \textbf{Requerimiento} & \textbf{Descripción} \\
\hline
\endhead
RF-05 & Crear reporte & El usuario puede crear reportes de incidentes con título, descripción, tipo, ubicación e imágenes \\
\hline
RF-06 & Categorización de reportes & Los reportes se clasifican por tipo: robo, incendio, emergencia, accidente, otro \\
\hline
RF-07 & Gestión de estados & Los reportes tienen estados: activo, en\_proceso, resuelto \\
\hline
RF-08 & Edición de reportes & El usuario puede editar sus reportes dentro de las primeras 24 horas \\
\hline
RF-09 & Listado de reportes & El sistema muestra todos los reportes activos con filtros por tipo y estado \\
\hline
RF-10 & Detalle de reporte & El usuario puede ver el detalle completo de un reporte incluyendo imágenes y ubicación en mapa \\
\hline
RF-11 & Adjuntar imágenes & El usuario puede adjuntar hasta 3 imágenes por reporte desde cámara o galería \\
\hline
\caption{Requerimientos funcionales - Reportes}
\end{longtable}

\subsection{Módulo de Chat Comunitario}

\begin{longtable}{|p{1.5cm}|p{4cm}|p{8cm}|}
\hline
\textbf{ID} & \textbf{Requerimiento} & \textbf{Descripción} \\
\hline
\endfirsthead
\hline
\textbf{ID} & \textbf{Requerimiento} & \textbf{Descripción} \\
\hline
\endhead
RF-12 & Chat grupal por zonas & El sistema provee grupos de chat predefinidos por zonas geográficas (Centro, Sur, Norte) \\
\hline
RF-13 & Crear grupos de chat & Los usuarios pueden crear grupos de chat públicos o privados \\
\hline
RF-14 & Enviar mensajes & Los usuarios pueden enviar mensajes de texto e imágenes en los grupos \\
\hline
RF-15 & Unirse/Salir de grupos & Los usuarios pueden unirse o abandonar grupos de chat \\
\hline
RF-16 & Chat privado & Los usuarios pueden iniciar conversaciones privadas entre dos personas \\
\hline
\caption{Requerimientos funcionales - Chat}
\end{longtable}

\subsection{Módulo de Administración}

\begin{longtable}{|p{1.5cm}|p{4cm}|p{8cm}|}
\hline
\textbf{ID} & \textbf{Requerimiento} & \textbf{Descripción} \\
\hline
\endfirsthead
\hline
\textbf{ID} & \textbf{Requerimiento} & \textbf{Descripción} \\
\hline
\endhead
RF-17 & Dashboard de estadísticas & Los administradores pueden ver estadísticas de reportes y actividad \\
\hline
RF-18 & Gestión de usuarios & Los administradores pueden ver, activar/desactivar usuarios \\
\hline
RF-19 & Moderación de reportes & Los administradores pueden cambiar el estado de cualquier reporte \\
\hline
RF-20 & Eliminación lógica & Los administradores pueden eliminar reportes (eliminación lógica) \\
\hline
\caption{Requerimientos funcionales - Administración}
\end{longtable}

\subsection{Módulo de Notificaciones}

\begin{longtable}{|p{1.5cm}|p{4cm}|p{8cm}|}
\hline
\textbf{ID} & \textbf{Requerimiento} & \textbf{Descripción} \\
\hline
\endfirsthead
\hline
\textbf{ID} & \textbf{Requerimiento} & \textbf{Descripción} \\
\hline
\endhead
RF-21 & Notificaciones de nuevos reportes & El sistema notifica a todos los usuarios cuando se crea un nuevo reporte \\
\hline
RF-22 & Notificaciones de cambio de estado & El sistema notifica al dueño cuando cambia el estado de su reporte \\
\hline
RF-23 & Notificaciones de chat & El sistema notifica cuando hay nuevos mensajes en los grupos \\
\hline
\caption{Requerimientos funcionales - Notificaciones}
\end{longtable}

\subsection{Módulo de Configuración}

\begin{longtable}{|p{1.5cm}|p{4cm}|p{8cm}|}
\hline
\textbf{ID} & \textbf{Requerimiento} & \textbf{Descripción} \\
\hline
\endfirsthead
\hline
\textbf{ID} & \textbf{Requerimiento} & \textbf{Descripción} \\
\hline
\endhead
RF-24 & Configuración de notificaciones & El usuario puede activar/desactivar notificaciones, sonido y vibración \\
\hline
RF-25 & Tema de la aplicación & El usuario puede elegir entre tema claro, oscuro o automático (sistema) \\
\hline
RF-26 & Limpiar caché & El usuario puede limpiar datos temporales de la aplicación \\
\hline
\caption{Requerimientos funcionales - Configuración}
\end{longtable}

% ============================================================================
% SECCIÓN 5: REQUERIMIENTOS NO FUNCIONALES
% ============================================================================

\section{Requerimientos No Funcionales (RNF)}

\subsection{Seguridad}

\begin{longtable}{|p{1.5cm}|p{4cm}|p{8cm}|}
\hline
\textbf{ID} & \textbf{Requerimiento} & \textbf{Implementación} \\
\hline
\endfirsthead
\hline
\textbf{ID} & \textbf{Requerimiento} & \textbf{Implementación} \\
\hline
\endhead
RNF-01 & Autenticación segura & Firebase Auth con email/contraseña \\
\hline
RNF-02 & Autorización por roles & Sistema de roles (user/admin) verificado en cliente y Firestore Rules \\
\hline
RNF-03 & Protección de datos sensibles & Enmascaramiento de ubicación exacta para usuarios no propietarios del reporte \\
\hline
RNF-04 & Validación de formularios & Validación de email, contraseña mínima 6 caracteres, campos requeridos \\
\hline
\caption{Requerimientos no funcionales - Seguridad}
\end{longtable}

\subsection{Rendimiento}

\begin{longtable}{|p{1.5cm}|p{4cm}|p{8cm}|}
\hline
\textbf{ID} & \textbf{Requerimiento} & \textbf{Implementación} \\
\hline
\endfirsthead
\hline
\textbf{ID} & \textbf{Requerimiento} & \textbf{Implementación} \\
\hline
\endhead
RNF-05 & Carga eficiente de datos & Uso de Streams de Firestore con límites (100 mensajes por carga) \\
\hline
RNF-06 & Optimización de imágenes & Compresión de imágenes antes de subir a Firebase Storage \\
\hline
RNF-07 & Caché local & SharedPreferences para configuraciones y preferencias \\
\hline
RNF-08 & Lazy loading & Carga bajo demanda de datos en listas \\
\hline
\caption{Requerimientos no funcionales - Rendimiento}
\end{longtable}

\subsection{Compatibilidad}

\begin{longtable}{|p{1.5cm}|p{4cm}|p{8cm}|}
\hline
\textbf{ID} & \textbf{Requerimiento} & \textbf{Implementación} \\
\hline
\endfirsthead
\hline
\textbf{ID} & \textbf{Requerimiento} & \textbf{Implementación} \\
\hline
\endhead
RNF-09 & Multiplataforma & Soporte para Android (principal), iOS y Web \\
\hline
RNF-10 & Responsive design & UI adaptable a diferentes tamaños de pantalla \\
\hline
RNF-11 & Soporte offline parcial & Firestore offline persistence habilitado \\
\hline
\caption{Requerimientos no funcionales - Compatibilidad}
\end{longtable}

\subsection{Usabilidad}

\begin{longtable}{|p{1.5cm}|p{4cm}|p{8cm}|}
\hline
\textbf{ID} & \textbf{Requerimiento} & \textbf{Implementación} \\
\hline
\endfirsthead
\hline
\textbf{ID} & \textbf{Requerimiento} & \textbf{Implementación} \\
\hline
\endhead
RNF-12 & Material Design 3 & Uso de ColorScheme.fromSeed para consistencia visual \\
\hline
RNF-13 & Feedback visual & Estados de carga, mensajes de error/éxito, indicadores de progreso \\
\hline
RNF-14 & Navegación intuitiva & Menú principal con acceso directo a funcionalidades según rol \\
\hline
RNF-15 & Accesibilidad & Uso de semántica de Flutter, tooltips en iconos \\
\hline
\caption{Requerimientos no funcionales - Usabilidad}
\end{longtable}

\subsection{Mantenibilidad}

\begin{longtable}{|p{1.5cm}|p{4cm}|p{8cm}|}
\hline
\textbf{ID} & \textbf{Requerimiento} & \textbf{Implementación} \\
\hline
\endfirsthead
\hline
\textbf{ID} & \textbf{Requerimiento} & \textbf{Implementación} \\
\hline
\endhead
RNF-16 & Código limpio & Separación de responsabilidades (MVVM), uso de flutter\_lints \\
\hline
RNF-17 & Modularidad & Servicios independientes con ChangeNotifier \\
\hline
RNF-18 & Documentación en código & Comentarios con referencias a RF (ej: // RF-05: Crear reporte) \\
\hline
\caption{Requerimientos no funcionales - Mantenibilidad}
\end{longtable}

% ============================================================================
% SECCIÓN 6: DIAGRAMAS
% ============================================================================

\section{Diagramas Técnicos}

\textbf{Nota:} Los diagramas se generan con Mermaid.js. Ver archivo separado \texttt{DIAGRAMAS\_MERMAID.md} para los códigos fuente.

\subsection{Diagrama de Casos de Uso}
Muestra la interacción entre los actores (Usuario y Administrador) con los requerimientos funcionales del sistema.

\subsection{Diagrama de Clases}
Representa la estructura de los modelos de datos (UserModel, Report, ChatGroup, Message) y los servicios (AuthService, ReportService, ChatService) con sus relaciones.

\subsection{Diagrama de Secuencia}
Ilustra el flujo de creación de un reporte, desde la interacción del usuario hasta el almacenamiento en Firebase y el envío de notificaciones.

% ============================================================================
% SECCIÓN 7: GUÍA DE INSTALACIÓN
% ============================================================================

\section{Guía de Instalación y Despliegue}

\subsection{Requisitos Previos}

\begin{table}[h]
\centering
\begin{tabular}{|l|l|l|}
\hline
\textbf{Herramienta} & \textbf{Versión Mínima} & \textbf{Verificar} \\
\hline
Flutter SDK & 3.7.2+ & \texttt{flutter --version} \\
Dart SDK & 3.7.2+ & \texttt{dart --version} \\
Android Studio & 2022.1+ & Con Android SDK 33+ \\
Git & 2.x & \texttt{git --version} \\
Cuenta Firebase & - & console.firebase.google.com \\
\hline
\end{tabular}
\caption{Requisitos previos}
\end{table}

\subsection{Pasos de Instalación}

\subsubsection{Paso 1: Clonar el Repositorio}
\begin{lstlisting}[language=bash]
git clone <URL_DEL_REPOSITORIO>
cd proyecto_safearea
\end{lstlisting}

\subsubsection{Paso 2: Instalar Dependencias}
\begin{lstlisting}[language=bash]
flutter pub get
\end{lstlisting}

\subsubsection{Paso 3: Configurar Firebase}
\begin{enumerate}
    \item Crear proyecto en Firebase Console
    \item Habilitar servicios: Authentication, Cloud Firestore, Storage, Cloud Messaging
    \item Descargar \texttt{google-services.json} → \texttt{android/app/}
    \item Descargar \texttt{GoogleService-Info.plist} → \texttt{ios/Runner/}
\end{enumerate}

\subsubsection{Paso 4: Configurar Google Maps}
\begin{enumerate}
    \item Obtener API Key en Google Cloud Console
    \item Habilitar "Maps SDK for Android"
    \item Agregar en \texttt{android/app/src/main/AndroidManifest.xml}:
\end{enumerate}

\begin{lstlisting}[language=xml]
<meta-data
    android:name="com.google.android.geo.API_KEY"
    android:value="TU_API_KEY"/>
\end{lstlisting}

\subsubsection{Paso 5: Ejecutar la Aplicación}
\begin{lstlisting}[language=bash]
# Listar dispositivos
flutter devices

# Ejecutar en dispositivo
flutter run -d <DEVICE_ID>

# Ejecutar en modo release
flutter run --release
\end{lstlisting}

\subsubsection{Paso 6: Generar APK de Producción}
\begin{lstlisting}[language=bash]
# APK release
flutter build apk --release

# Bundle para Play Store
flutter build appbundle --release
\end{lstlisting}

\subsection{Índices de Firestore Requeridos}

\begin{table}[h]
\centering
\begin{tabular}{|l|l|}
\hline
\textbf{Colección} & \textbf{Campos} \\
\hline
reports & isActive (ASC), createdAt (DESC) \\
reports & isActive (ASC), type (ASC), createdAt (DESC) \\
reports & userId (ASC), isActive (ASC), createdAt (DESC) \\
chatGroups & isPublic (ASC), createdAt (DESC) \\
\hline
\end{tabular}
\caption{Índices compuestos de Firestore}
\end{table}

% ============================================================================
% SECCIÓN 8: MANUAL DE USUARIO
% ============================================================================

\section{Manual de Usuario}

\subsection{Primeros Pasos}

\subsubsection{Registro de Cuenta}
\begin{enumerate}
    \item Abre la aplicación SafeArea
    \item Presiona "Regístrate aquí"
    \item Completa: Nombre, Correo electrónico, Contraseña (mín. 6 caracteres)
    \item Presiona "Registrarse"
\end{enumerate}

\subsubsection{Inicio de Sesión}
\begin{enumerate}
    \item Ingresa tu correo electrónico
    \item Ingresa tu contraseña
    \item Presiona "Iniciar Sesión"
\end{enumerate}

\subsection{Pantalla Principal (Home)}

\begin{table}[h]
\centering
\begin{tabular}{|l|l|}
\hline
\textbf{Opción} & \textbf{Descripción} \\
\hline
Reportes & Ver y crear reportes de incidentes \\
Chat & Acceder a grupos de chat por zonas \\
Configuración & Ajustes de la aplicación \\
Mi Perfil & Ver y editar información personal \\
\hline
\end{tabular}
\caption{Opciones del menú principal}
\end{table}

\textit{Los administradores verán opciones adicionales: Dashboard y Usuarios}

\subsection{Crear un Reporte}

\begin{enumerate}
    \item Desde Home, presiona "Reportes"
    \item Presiona el botón "+" (FAB azul)
    \item Completa el formulario:
    \begin{itemize}
        \item Tipo de incidente: Robo, Incendio, Emergencia, Accidente, Otro
        \item Título: Descripción breve del incidente
        \item Descripción: Detalles del incidente
        \item Ubicación: Selecciona en el mapa
        \item Imágenes (opcional): Adjunta hasta 3 fotos
    \end{itemize}
    \item Presiona "Crear Reporte"
\end{enumerate}

\subsection{Usar el Chat Comunitario}

\begin{enumerate}
    \item Desde Home, presiona "Chat"
    \item Selecciona una zona: Centro, Sur, Norte
    \item Escribe tu mensaje y presiona enviar
    \item Puedes adjuntar imágenes con el icono de cámara
\end{enumerate}

\subsection{Configurar la Aplicación}

\begin{enumerate}
    \item Desde Home, presiona "Configuración"
    \item Opciones disponibles:
    \begin{itemize}
        \item Notificaciones: Activar/desactivar alertas
        \item Apariencia: Tema claro, oscuro o automático
        \item Información: Versión de la app
        \item Datos: Limpiar caché
    \end{itemize}
\end{enumerate}

% ============================================================================
% SECCIÓN 9: INFORMACIÓN DEL PROYECTO
% ============================================================================

\section{Información del Proyecto}

\begin{table}[h]
\centering
\begin{tabular}{|l|l|}
\hline
\textbf{Campo} & \textbf{Valor} \\
\hline
Nombre & SafeArea \\
Versión & 0.1.0 \\
Package & com.newsdevs.safearea \\
Plataformas & Android, iOS, Web \\
Idioma & Español \\
\hline
\end{tabular}
\caption{Información del proyecto}
\end{table}

\vfill

\begin{center}
\textit{Documentación generada el 27 de Noviembre de 2025}
\end{center}

\end{document}
